\section{Theoretical Background}
\label{sec:background}

This section presents the theoretical studies that support the experience reported.

\subsection{Heuristic evaluation and the pedagogical challenge}

Heuristic evaluation is an inspection-based usability evaluation method introduced by Nielsen and Molich \cite{nielsen1990heuristic}, in which evaluators examine an interface against a small set of usability principles and report the problems they find. The most widely used set is the ten heuristics consolidated by Nielsen \cite{nielsen1994heuristic}, with phrasing periodically refined by the Nielsen Norman Group \cite{nielsen2024tenheuristics} and reproduced in Brazilian HCI textbooks \cite{barbosa2010interacao, BarbosaEtAl_Livro2021}. Compared with empirical user testing, heuristic evaluation is cheap and fast to apply, which is one of the reasons for its wide use.

Recognizing violations of the heuristics in real interfaces, however, is not trivial. Nielsen \cite{nielsen1992finding} reports a clear gap: novice evaluators identify substantially fewer usability problems than regular usability specialists, and ``double experts'' --- evaluators with both usability and application-domain expertise --- perform better still. A wide evaluator effect further shows substantial variability between evaluators applying the same set of heuristics \cite{hertzum2003evaluator}, and training has been identified as a key factor influencing inspection outcomes \cite{hvannberg2007heuristic}. More recent work points to abstract heuristic descriptions as a specific obstacle for novices and proposes rephrasing the heuristics in more concrete language to help comprehension \cite{abulfaraj2020detailed}.

Taken together, these outcomes make a practical case for HCI educators: reading the ten heuristics in isolation rarely provides learners with sufficient anchoring to identify the same ideas in their own evaluations. Concrete examples in familiar tasks are needed before novices can apply the heuristics confidently, with room for new teaching approaches.

\subsection{Active and game-based learning in computing education}

Active learning refers to pedagogical approaches, such as discussion, problem-solving, role-play, or guided practice, that engage learners in doing something with the material rather than passively receiving information \cite{bonwell1991active}. A large meta-analysis across STEM (Science, Technology, Engineering, and Mathematics) disciplines reports measurable gains in student performance under active instruction compared with traditional lectures \cite{freeman2014active}. Kolb's experiential learning cycle \cite{kolb1984experiential} --- concrete experience followed by reflection, conceptualization, and renewed experimentation --- is a frequent frame for designing such activities, since it links a hands-on encounter with an artifact to the abstract reasoning that comes after.

In computing education, and in HCI specifically, active practices take the form of design exercises, simulated evaluations, peer reviews, and games or playful artifacts that give learners concrete experience of the concepts under study \cite{silveira2024praticas}. In Brazil, HCI is taught across different computing curricula with varying depth and time allocation \cite{boscarioli2013hci,boscarioli2014charting}. Time-efficient classroom activities that fit a single session and require no specific infrastructure are therefore especially valuable in this teaching context.
